% P10-4c-Nachtrag, 15.09.2026: Modellwahl aus Iteration 8 sowie L3-B/L4-B
% (01.07.) und MC0-Stress (07.07.) eingeordnet; Belegdetails im P10-4c-Bericht.
% M16, 11.09.2026: Leserführung und Absatzlogik; Anforderungen und Anker erhalten.
% ============================================================
% §4.2 Verdichteter Explorationsverlauf — v05 (Claude)
% v05 (01.09., Kollegen-Runde 4.2 — kritisch geprüft): ① Zählfehler
% „neun" → „Z0 bis Z9" (waren zehn; mein Fehler) ② Synthese-Absatz:
% Verantwortlichkeiten-Satz sprachlich + „menschlich autorisiert"
% ergänzt (Vierer-Teilung Agentik/Workflow/Deterministik/Mensch)
% ③ Modellhistorie NICHT abgeschwächt — Kausalität ist Iter-8-belegt
% (Kollegen-Bedingung erfüllt) ④ Fixierungs-Satz geglättet
% ⑤ Vergleichsarm konkretisiert (A' = freie Variante unter denselben
% Modellbedingungen; verifiziert am Spike) ⑥ Schluss direkter
% („ausgewählte Wirkmechanismen … Kapitel-7-Bedingungen"); Tabelle:
% Z6-Designrichtung ohne „Maker–Checker–Repair" (Begriff wird erst in
% 4.5 geprägt — roter Faden).
% v04 (01.09., Konkretisierungs-Runde): M1 Modellhistorie-Satz im
% P0-Absatz (lokal → gehostet → fester Arbeitsstandard; defensiv per
% E4-Entscheid, Toolcall-Ausfälle notes-belegt, keine Modellwertung);
% Autor-X-Kommentare verarbeitet+entfernt.
% v03 (01.09., „dosierte Agentik"-Runde): Synthese-Absatz nach der
% Tabelle — der übergeordnete Entwicklungspfad (Differenzierung der
% Verantwortlichkeiten, NICHT Ablösung von Agenten durch Workflows).
% v02 (01.09., Kollegen-Review): „ordnet…zu" statt „Inhaltsverzeichnis";
% Z9-Schlusssatz expliziert (Stabilisierung → Kap. 6 Realisierung /
% Kap. 7 Evaluation); P0-Absatz präzisiert (Fixierung INNERHALB der
% Vergleichsanordnung; „keine kausale Isolation eines Designeffekts").
%
% Grundlage: Kap-4-Blaupause §7 (Tabelle statt Chronik-Prosa;
% P0 = 2–3 Sätze nach der Tabelle mit VORGEGEBENER zulässiger
% Aussage; Grenz-Absatz) · genese-dossier §2 (Z0–Z9, auditiert
% inkl. Korrektur-Kopf) · Kap-3-v02 §13.2 (KEINE Wiederholung der
% Vier-Phasen aus 3.2 — nur inhaltliche Orientierung).
%
% Entscheidungen:
%   · Die Anker-Spalte der Tabelle verweist auf E-IDs + Tags —
%     der Zyklen-Überblick wird damit zum Inhaltsverzeichnis des
%     Evidenzsystems (statt eigener Beleg-Doppelung).
%   · Z3 trägt die korrigierte Wahrheit („zwei von drei geplanten
%     Varianten"); Z7 deklariert die Zielbild-Herkunft direkt in
%     der Tabelle („kein Fehlschlag-Pivot").
%   · P0-Absatz = Blaupausen-Wortlaut („explorativer Vergleichsarm",
%     „reduzierten den Modelleffekt als alternative Erklärung",
%     KEINE kausale Isolation).
% Budget ~0,9 S. inkl. Tabelle (Ziel 0,8–1,0).
% [Fundstellen P0: genese-dossier §3/P0 + thesis-evidence/README
%  Spike 5 (A'-Arm); Modell-Wechsel: phase01-zsmfassung (Ollama),
%  Run-Configs Ende Mai (gpt-oss-120b:free / gpt-4.1-mini)]
% ============================================================

\section{Verdichteter Explorationsverlauf}
\label{sec:explorationsverlauf}

Tabelle~\ref{t:explorationsverlauf-zyklen} ordnet die prägenden
Entwicklungsschritte in die Zyklen Z0 bis Z9. Sie verbindet den jeweiligen
Arbeitsfokus mit einer Beobachtung oder Zielkonkretisierung, der daraus
abgeleiteten Designrichtung und den zugehörigen Belegen. Die Zyklen
strukturieren das in Abschnitt~\ref{sec:iteratives-vorgehen-evidenz}
beschriebene Vorgehen rückblickend. Sie bilden keine strikt getrennten
Zeiträume ab; einzelne Zwischenlösungen bestanden parallel oder wurden
verworfen.

Die folgenden Abschnitte vertiefen diese Entwicklung: Z0 bis Z3 betreffen
Ausführung und Informationsführung
(Abschnitt~\ref{sec:beobachtbare-ausfuehrung-informationsfuehrung}),
Z4 und Z5 die Quellenbindung
(Abschnitt~\ref{sec:evidenzgebundene-wissensgewinnung}) und Z6 die
Trennung der Verarbeitungsverantwortlichkeiten
(Abschnitt~\ref{sec:kontrollierte-semantische-transformation}). Z7
verbindet Projektzustand und Fortschreibung
(Abschnitte~\ref{sec:projektzustand-kontrollierte-mutation}
und~\ref{sec:einordnende-fortschreibung}); Z8 ergänzt Außenkopplung
und kontinuierliche Nutzung
(Abschnitte~\ref{sec:projektionen-aussenkopplung}
und~\ref{sec:durabilitaet-kontinuierliche-nutzung}). Z9 bezeichnet die
Stabilisierung und technische Prüfung. Die Realisierung und die
untersuchten Lauf- und Teststände beschreiben
Kapitel~\ref{chap:maf-realisierung} und~\ref{chap:evaluation}.

% P10-4c, 14.09.2026: Ledger-Herleitung und Nachweisanschlüsse präzisiert.
% B-46-Nachtrag, 13.09.2026: Z3 an die Originalprüfung E43-05
% und die bereits synchronisierten Fassungen von 4.3/Anhang A angeglichen.
% M16 / B-34, 11.09.2026: Z4 trennt die Artefakte der beiden Prüfbeobachtungen.
% M03 / B-17, 10.09.2026: ursprüngliches und erweitertes Nutzungsziel
% unterschieden; keine neue Datierung des Betreuerimpulses.
% ============================================================
% Tabelle zu §4.2 — Ziel gemäß aktueller Einbindung in chap4.2.tex:
%   tables/chap4/chap4.2table.tex
% Aufruf: % P10-4c, 14.09.2026: Ledger-Herleitung und Nachweisanschlüsse präzisiert.
% B-46-Nachtrag, 13.09.2026: Z3 an die Originalprüfung E43-05
% und die bereits synchronisierten Fassungen von 4.3/Anhang A angeglichen.
% M16 / B-34, 11.09.2026: Z4 trennt die Artefakte der beiden Prüfbeobachtungen.
% M03 / B-17, 10.09.2026: ursprüngliches und erweitertes Nutzungsziel
% unterschieden; keine neue Datierung des Betreuerimpulses.
% ============================================================
% Tabelle zu §4.2 — Ziel gemäß aktueller Einbindung in chap4.2.tex:
%   tables/chap4/chap4.2table.tex
% Aufruf: % P10-4c, 14.09.2026: Ledger-Herleitung und Nachweisanschlüsse präzisiert.
% B-46-Nachtrag, 13.09.2026: Z3 an die Originalprüfung E43-05
% und die bereits synchronisierten Fassungen von 4.3/Anhang A angeglichen.
% M16 / B-34, 11.09.2026: Z4 trennt die Artefakte der beiden Prüfbeobachtungen.
% M03 / B-17, 10.09.2026: ursprüngliches und erweitertes Nutzungsziel
% unterschieden; keine neue Datierung des Betreuerimpulses.
% ============================================================
% Tabelle zu §4.2 — Ziel gemäß aktueller Einbindung in chap4.2.tex:
%   tables/chap4/chap4.2table.tex
% Aufruf: \input{tables/chap4/chap4.2table}
% Label: t:explorationsverlauf-zyklen
% v03 (01.09., Tabellen-Feedback-Runde — kritisch geprüft): Z1 konkreter
% (Übergaben + Ablaufreihenfolge explizit zu gestalten) · Z2 „als eigene
% Designvariable" · Z4-Designrichtung „an Quellenrepräsentation BINDEN"
% (statt „verlagern") · Z6 ohne MCR-Begriff (fällt erst in 4.5) · Z7
% „deterministische Zuordnungslogik" (VERIFIZIERT am Audit: IngestionApply
% schrieb featureKey-String als Relationsziel — stärkere Variante belegt;
% 4.6-Text + E46-01-Zeile synchronisiert) · Z8 „im Entwicklungsbetrieb"
% (4.7-Grenze schon hier) + „unter Nutzung der bestehenden Governance"
% (= A8-Wortlaut) · Z9 als Stabilisierungs-Übergang · Caption
% „zugehörigen Ankern" (E-IDs+Tags+Kapitelverweise gemischt zulässig).
% v02 (01.09., Kollegen-Review): Zeilen Z0–Z6/Z8/Z9 auf Evidenzgrenzen
% nachgeschärft (keine unbelegten Wertungen/Kausalitäten mehr: Z0 ohne
% „Ergebnisqualität modellabhängig", Z1 als Zielkonkretisierung, Z3
% „unbelegt trotz geschärfter Promptvorgaben", Z4 ohne Zerlegungs-
% Kausalität, Z5 Zirkularität als Eignungsaussage, Z6 Variabilitäts-
% beschreibung statt „inhaltlich scharf", Z8 „regulärer Ablaufzustand",
% Z9 „technische Tests").
% Grundlage: genese-dossier §2 (Z0–Z9), AUDITIERT 31.08. inkl.
% Korrektur-Kopf (A/B/C: zwei ausgeführt, C nie implementiert).
% Systemnamen vermieden; Evidenzanker-Spalte verweist auf die E-IDs
% aus Anhang A und die Artefaktstände (git-Tags) — verbindet Zyklen-
% Überblick und Evidenzsystem.
% ============================================================

\begin{table}[tbp]
\centering
\caption[Verdichteter Explorationsverlauf]{Verdichteter
Explorationsverlauf: Entwicklungszyklen mit prägender Beobachtung
beziehungsweise Zielkonkretisierung, daraus folgender Designrichtung und
den zugehörigen Ankern (E-IDs nach Anhang~\ref{anh:evidenzanker},
markierte Artefaktstände \texttt{v-s}$n$, Kapitelverweise). Eigene
Darstellung.}
\label{t:explorationsverlauf-zyklen}
\scriptsize
\setlength{\tabcolsep}{3pt}
\renewcommand{\arraystretch}{1.2}
\begin{tabularx}{\textwidth}{
    >{\raggedright\arraybackslash}p{0.03\textwidth}
    >{\raggedright\arraybackslash}p{0.16\textwidth}
    >{\raggedright\arraybackslash}X
    >{\raggedright\arraybackslash}p{0.20\textwidth}
    >{\raggedright\arraybackslash}p{0.115\textwidth}
}
\toprule
& \textbf{Fokus} &
\textbf{Prägende Beobachtung / Zielkonkretisierung} &
\textbf{Designrichtung} & \textbf{Anker} \\
\midrule
Z0 & Einzelagent-Durchstich &
Modelltext beschrieb Werkzeugaufrufe ohne registrierte Ausführung &
Ausführung unabhängig von der Modellausgabe beobachtbar machen &
E43-01, E43-02; \texttt{v-s0} \\
\addlinespace[0.3em]
Z1 & linearer Mehr-Agenten-Ablauf &
mit der Aufteilung auf mehrere spezialisierte Rollen mussten auch deren
Informationsübergaben und Ablaufreihenfolge explizit gestaltet werden &
Informationsübergaben und Ablaufsteuerung getrennt von den fachlichen
Rollen betrachten &
\texttt{v-s1} \\
\addlinespace[0.3em]
Z2 & Beobachtbarkeit ausgebaut &
Artefakterzeugung ohne vorausgehenden eigenen Quellzugriff; tatsächlicher
Informationsweg auf Verarbeitungsebene nicht explizit ausgewiesen &
Kontextführung als eigene Designvariable untersuchen &
E43-03; \texttt{v-s2} \\
\addlinespace[0.3em]
Z3 & Strategievergleich (zwei von drei geplanten Varianten) &
nach Promptschärfung zunächst Schreiben ohne Leseaufruf,
im nächsten Lauf Lesen ohne Artefakterzeugung; zwei Varianten
der Kontextführung wurden explorativ verglichen &
Informationsflüsse explizit modellieren &
E43-04, E43-05 \\
\addlinespace[0.3em]
Z4 & nachgelagerte Prüfung als Messinstrument &
freie Prüfung variierte und übersah ausgewählte Inhaltslücken;
festgelegte Prüfeinheiten und Themen ermöglichten gezielte
Einzelprüfungen; die Vollständigkeitsfrage blieb offen &
Vollständigkeitsprüfung an eine explizite Quellenrepräsentation binden &
E44-01, E44-04; \texttt{v-s3} \\
\addlinespace[0.3em]
Z5 & quellgebundene Zwischenrepräsentation &
keine explizite Zuordnung Artefakt--Quellstelle; eine allein
modellgenerierte, nicht menschlich bestätigte Vollständigkeitsreferenz
war als unabhängiger
Bewertungsmaßstab ungeeignet &
Evidenz vor Erzeugung; geteilte Referenzbildung mit menschlicher
Bestätigung &
E44-02, E44-03, E44-05; \texttt{v-s4} \\
\addlinespace[0.3em]
Z6 & getrennte Prüf- und Erzeugungsverantwortung &
freie Prüfung variierte zwischen Wiederholungen in Zahl, Kategorien und
Granularität der Befunde &
Erzeugung, Prüfung und Reparatur als getrennte
Workflow-Verantwortlichkeiten &
E45-02, E45-03; \texttt{v-s5}--\texttt{s6} \\
\addlinespace[0.3em]
Z7 & laufübergreifender Zustand, Kontrolle und Fortschreibung &
Konkretisierung des erweiterten Nutzungsziels (kein Fehlschlag-Pivot);
später: fehlerhafte
Relationen aus einer deterministisch implementierten Zuordnungsregel &
Zustand mit Identitäten, Relationen, Historie; gemeinsamer Mutationspfad
mit Invariantenprüfung; Einordnung neuer Information gegen den Bestand &
E46-01, E47-02; \texttt{v-s7}--\texttt{s8} \\
\addlinespace[0.3em]
Z8 & Außenkopplung, durable Ausführung und Bedienschicht &
menschliche Entscheidungspunkte machen planmäßiges Anhalten zu einem
regulären Ablaufzustand; direkte Einzelbedienung wurde im
Entwicklungsbetrieb mit wachsender Systembreite aufwendiger &
Projektionen aus dem Zustand mit kontrolliertem Rückweg;
Sicherungspunkte und Wiederaufnahme; agentische Bedienebene unter
Nutzung der bestehenden Governance &
E47-01, E48-01/-02; \texttt{v-s9} \\
\addlinespace[0.3em]
Z9 & Stabilisierung des Untersuchungsartefakts &
End-to-End-Abnahmen, Integritätsprüfungen und technische Tests &
stabilisierten Artefaktstand für die Evaluation herstellen &
Kapitel~\ref{chap:maf-realisierung}/\ref{chap:evaluation} \\
\bottomrule
\end{tabularx}
\end{table}

% Label: t:explorationsverlauf-zyklen
% v03 (01.09., Tabellen-Feedback-Runde — kritisch geprüft): Z1 konkreter
% (Übergaben + Ablaufreihenfolge explizit zu gestalten) · Z2 „als eigene
% Designvariable" · Z4-Designrichtung „an Quellenrepräsentation BINDEN"
% (statt „verlagern") · Z6 ohne MCR-Begriff (fällt erst in 4.5) · Z7
% „deterministische Zuordnungslogik" (VERIFIZIERT am Audit: IngestionApply
% schrieb featureKey-String als Relationsziel — stärkere Variante belegt;
% 4.6-Text + E46-01-Zeile synchronisiert) · Z8 „im Entwicklungsbetrieb"
% (4.7-Grenze schon hier) + „unter Nutzung der bestehenden Governance"
% (= A8-Wortlaut) · Z9 als Stabilisierungs-Übergang · Caption
% „zugehörigen Ankern" (E-IDs+Tags+Kapitelverweise gemischt zulässig).
% v02 (01.09., Kollegen-Review): Zeilen Z0–Z6/Z8/Z9 auf Evidenzgrenzen
% nachgeschärft (keine unbelegten Wertungen/Kausalitäten mehr: Z0 ohne
% „Ergebnisqualität modellabhängig", Z1 als Zielkonkretisierung, Z3
% „unbelegt trotz geschärfter Promptvorgaben", Z4 ohne Zerlegungs-
% Kausalität, Z5 Zirkularität als Eignungsaussage, Z6 Variabilitäts-
% beschreibung statt „inhaltlich scharf", Z8 „regulärer Ablaufzustand",
% Z9 „technische Tests").
% Grundlage: genese-dossier §2 (Z0–Z9), AUDITIERT 31.08. inkl.
% Korrektur-Kopf (A/B/C: zwei ausgeführt, C nie implementiert).
% Systemnamen vermieden; Evidenzanker-Spalte verweist auf die E-IDs
% aus Anhang A und die Artefaktstände (git-Tags) — verbindet Zyklen-
% Überblick und Evidenzsystem.
% ============================================================

\begin{table}[tbp]
\centering
\caption[Verdichteter Explorationsverlauf]{Verdichteter
Explorationsverlauf: Entwicklungszyklen mit prägender Beobachtung
beziehungsweise Zielkonkretisierung, daraus folgender Designrichtung und
den zugehörigen Ankern (E-IDs nach Anhang~\ref{anh:evidenzanker},
markierte Artefaktstände \texttt{v-s}$n$, Kapitelverweise). Eigene
Darstellung.}
\label{t:explorationsverlauf-zyklen}
\scriptsize
\setlength{\tabcolsep}{3pt}
\renewcommand{\arraystretch}{1.2}
\begin{tabularx}{\textwidth}{
    >{\raggedright\arraybackslash}p{0.03\textwidth}
    >{\raggedright\arraybackslash}p{0.16\textwidth}
    >{\raggedright\arraybackslash}X
    >{\raggedright\arraybackslash}p{0.20\textwidth}
    >{\raggedright\arraybackslash}p{0.115\textwidth}
}
\toprule
& \textbf{Fokus} &
\textbf{Prägende Beobachtung / Zielkonkretisierung} &
\textbf{Designrichtung} & \textbf{Anker} \\
\midrule
Z0 & Einzelagent-Durchstich &
Modelltext beschrieb Werkzeugaufrufe ohne registrierte Ausführung &
Ausführung unabhängig von der Modellausgabe beobachtbar machen &
E43-01, E43-02; \texttt{v-s0} \\
\addlinespace[0.3em]
Z1 & linearer Mehr-Agenten-Ablauf &
mit der Aufteilung auf mehrere spezialisierte Rollen mussten auch deren
Informationsübergaben und Ablaufreihenfolge explizit gestaltet werden &
Informationsübergaben und Ablaufsteuerung getrennt von den fachlichen
Rollen betrachten &
\texttt{v-s1} \\
\addlinespace[0.3em]
Z2 & Beobachtbarkeit ausgebaut &
Artefakterzeugung ohne vorausgehenden eigenen Quellzugriff; tatsächlicher
Informationsweg auf Verarbeitungsebene nicht explizit ausgewiesen &
Kontextführung als eigene Designvariable untersuchen &
E43-03; \texttt{v-s2} \\
\addlinespace[0.3em]
Z3 & Strategievergleich (zwei von drei geplanten Varianten) &
nach Promptschärfung zunächst Schreiben ohne Leseaufruf,
im nächsten Lauf Lesen ohne Artefakterzeugung; zwei Varianten
der Kontextführung wurden explorativ verglichen &
Informationsflüsse explizit modellieren &
E43-04, E43-05 \\
\addlinespace[0.3em]
Z4 & nachgelagerte Prüfung als Messinstrument &
freie Prüfung variierte und übersah ausgewählte Inhaltslücken;
festgelegte Prüfeinheiten und Themen ermöglichten gezielte
Einzelprüfungen; die Vollständigkeitsfrage blieb offen &
Vollständigkeitsprüfung an eine explizite Quellenrepräsentation binden &
E44-01, E44-04; \texttt{v-s3} \\
\addlinespace[0.3em]
Z5 & quellgebundene Zwischenrepräsentation &
keine explizite Zuordnung Artefakt--Quellstelle; eine allein
modellgenerierte, nicht menschlich bestätigte Vollständigkeitsreferenz
war als unabhängiger
Bewertungsmaßstab ungeeignet &
Evidenz vor Erzeugung; geteilte Referenzbildung mit menschlicher
Bestätigung &
E44-02, E44-03, E44-05; \texttt{v-s4} \\
\addlinespace[0.3em]
Z6 & getrennte Prüf- und Erzeugungsverantwortung &
freie Prüfung variierte zwischen Wiederholungen in Zahl, Kategorien und
Granularität der Befunde &
Erzeugung, Prüfung und Reparatur als getrennte
Workflow-Verantwortlichkeiten &
E45-02, E45-03; \texttt{v-s5}--\texttt{s6} \\
\addlinespace[0.3em]
Z7 & laufübergreifender Zustand, Kontrolle und Fortschreibung &
Konkretisierung des erweiterten Nutzungsziels (kein Fehlschlag-Pivot);
später: fehlerhafte
Relationen aus einer deterministisch implementierten Zuordnungsregel &
Zustand mit Identitäten, Relationen, Historie; gemeinsamer Mutationspfad
mit Invariantenprüfung; Einordnung neuer Information gegen den Bestand &
E46-01, E47-02; \texttt{v-s7}--\texttt{s8} \\
\addlinespace[0.3em]
Z8 & Außenkopplung, durable Ausführung und Bedienschicht &
menschliche Entscheidungspunkte machen planmäßiges Anhalten zu einem
regulären Ablaufzustand; direkte Einzelbedienung wurde im
Entwicklungsbetrieb mit wachsender Systembreite aufwendiger &
Projektionen aus dem Zustand mit kontrolliertem Rückweg;
Sicherungspunkte und Wiederaufnahme; agentische Bedienebene unter
Nutzung der bestehenden Governance &
E47-01, E48-01/-02; \texttt{v-s9} \\
\addlinespace[0.3em]
Z9 & Stabilisierung des Untersuchungsartefakts &
End-to-End-Abnahmen, Integritätsprüfungen und technische Tests &
stabilisierten Artefaktstand für die Evaluation herstellen &
Kapitel~\ref{chap:maf-realisierung}/\ref{chap:evaluation} \\
\bottomrule
\end{tabularx}
\end{table}

% Label: t:explorationsverlauf-zyklen
% v03 (01.09., Tabellen-Feedback-Runde — kritisch geprüft): Z1 konkreter
% (Übergaben + Ablaufreihenfolge explizit zu gestalten) · Z2 „als eigene
% Designvariable" · Z4-Designrichtung „an Quellenrepräsentation BINDEN"
% (statt „verlagern") · Z6 ohne MCR-Begriff (fällt erst in 4.5) · Z7
% „deterministische Zuordnungslogik" (VERIFIZIERT am Audit: IngestionApply
% schrieb featureKey-String als Relationsziel — stärkere Variante belegt;
% 4.6-Text + E46-01-Zeile synchronisiert) · Z8 „im Entwicklungsbetrieb"
% (4.7-Grenze schon hier) + „unter Nutzung der bestehenden Governance"
% (= A8-Wortlaut) · Z9 als Stabilisierungs-Übergang · Caption
% „zugehörigen Ankern" (E-IDs+Tags+Kapitelverweise gemischt zulässig).
% v02 (01.09., Kollegen-Review): Zeilen Z0–Z6/Z8/Z9 auf Evidenzgrenzen
% nachgeschärft (keine unbelegten Wertungen/Kausalitäten mehr: Z0 ohne
% „Ergebnisqualität modellabhängig", Z1 als Zielkonkretisierung, Z3
% „unbelegt trotz geschärfter Promptvorgaben", Z4 ohne Zerlegungs-
% Kausalität, Z5 Zirkularität als Eignungsaussage, Z6 Variabilitäts-
% beschreibung statt „inhaltlich scharf", Z8 „regulärer Ablaufzustand",
% Z9 „technische Tests").
% Grundlage: genese-dossier §2 (Z0–Z9), AUDITIERT 31.08. inkl.
% Korrektur-Kopf (A/B/C: zwei ausgeführt, C nie implementiert).
% Systemnamen vermieden; Evidenzanker-Spalte verweist auf die E-IDs
% aus Anhang A und die Artefaktstände (git-Tags) — verbindet Zyklen-
% Überblick und Evidenzsystem.
% ============================================================

\begin{table}[tbp]
\centering
\caption[Verdichteter Explorationsverlauf]{Verdichteter
Explorationsverlauf: Entwicklungszyklen mit prägender Beobachtung
beziehungsweise Zielkonkretisierung, daraus folgender Designrichtung und
den zugehörigen Ankern (E-IDs nach Anhang~\ref{anh:evidenzanker},
markierte Artefaktstände \texttt{v-s}$n$, Kapitelverweise). Eigene
Darstellung.}
\label{t:explorationsverlauf-zyklen}
\scriptsize
\setlength{\tabcolsep}{3pt}
\renewcommand{\arraystretch}{1.2}
\begin{tabularx}{\textwidth}{
    >{\raggedright\arraybackslash}p{0.03\textwidth}
    >{\raggedright\arraybackslash}p{0.16\textwidth}
    >{\raggedright\arraybackslash}X
    >{\raggedright\arraybackslash}p{0.20\textwidth}
    >{\raggedright\arraybackslash}p{0.115\textwidth}
}
\toprule
& \textbf{Fokus} &
\textbf{Prägende Beobachtung / Zielkonkretisierung} &
\textbf{Designrichtung} & \textbf{Anker} \\
\midrule
Z0 & Einzelagent-Durchstich &
Modelltext beschrieb Werkzeugaufrufe ohne registrierte Ausführung &
Ausführung unabhängig von der Modellausgabe beobachtbar machen &
E43-01, E43-02; \texttt{v-s0} \\
\addlinespace[0.3em]
Z1 & linearer Mehr-Agenten-Ablauf &
mit der Aufteilung auf mehrere spezialisierte Rollen mussten auch deren
Informationsübergaben und Ablaufreihenfolge explizit gestaltet werden &
Informationsübergaben und Ablaufsteuerung getrennt von den fachlichen
Rollen betrachten &
\texttt{v-s1} \\
\addlinespace[0.3em]
Z2 & Beobachtbarkeit ausgebaut &
Artefakterzeugung ohne vorausgehenden eigenen Quellzugriff; tatsächlicher
Informationsweg auf Verarbeitungsebene nicht explizit ausgewiesen &
Kontextführung als eigene Designvariable untersuchen &
E43-03; \texttt{v-s2} \\
\addlinespace[0.3em]
Z3 & Strategievergleich (zwei von drei geplanten Varianten) &
nach Promptschärfung zunächst Schreiben ohne Leseaufruf,
im nächsten Lauf Lesen ohne Artefakterzeugung; zwei Varianten
der Kontextführung wurden explorativ verglichen &
Informationsflüsse explizit modellieren &
E43-04, E43-05 \\
\addlinespace[0.3em]
Z4 & nachgelagerte Prüfung als Messinstrument &
freie Prüfung variierte und übersah ausgewählte Inhaltslücken;
festgelegte Prüfeinheiten und Themen ermöglichten gezielte
Einzelprüfungen; die Vollständigkeitsfrage blieb offen &
Vollständigkeitsprüfung an eine explizite Quellenrepräsentation binden &
E44-01, E44-04; \texttt{v-s3} \\
\addlinespace[0.3em]
Z5 & quellgebundene Zwischenrepräsentation &
keine explizite Zuordnung Artefakt--Quellstelle; eine allein
modellgenerierte, nicht menschlich bestätigte Vollständigkeitsreferenz
war als unabhängiger
Bewertungsmaßstab ungeeignet &
Evidenz vor Erzeugung; geteilte Referenzbildung mit menschlicher
Bestätigung &
E44-02, E44-03, E44-05; \texttt{v-s4} \\
\addlinespace[0.3em]
Z6 & getrennte Prüf- und Erzeugungsverantwortung &
freie Prüfung variierte zwischen Wiederholungen in Zahl, Kategorien und
Granularität der Befunde &
Erzeugung, Prüfung und Reparatur als getrennte
Workflow-Verantwortlichkeiten &
E45-02, E45-03; \texttt{v-s5}--\texttt{s6} \\
\addlinespace[0.3em]
Z7 & laufübergreifender Zustand, Kontrolle und Fortschreibung &
Konkretisierung des erweiterten Nutzungsziels (kein Fehlschlag-Pivot);
später: fehlerhafte
Relationen aus einer deterministisch implementierten Zuordnungsregel &
Zustand mit Identitäten, Relationen, Historie; gemeinsamer Mutationspfad
mit Invariantenprüfung; Einordnung neuer Information gegen den Bestand &
E46-01, E47-02; \texttt{v-s7}--\texttt{s8} \\
\addlinespace[0.3em]
Z8 & Außenkopplung, durable Ausführung und Bedienschicht &
menschliche Entscheidungspunkte machen planmäßiges Anhalten zu einem
regulären Ablaufzustand; direkte Einzelbedienung wurde im
Entwicklungsbetrieb mit wachsender Systembreite aufwendiger &
Projektionen aus dem Zustand mit kontrolliertem Rückweg;
Sicherungspunkte und Wiederaufnahme; agentische Bedienebene unter
Nutzung der bestehenden Governance &
E47-01, E48-01/-02; \texttt{v-s9} \\
\addlinespace[0.3em]
Z9 & Stabilisierung des Untersuchungsartefakts &
End-to-End-Abnahmen, Integritätsprüfungen und technische Tests &
stabilisierten Artefaktstand für die Evaluation herstellen &
Kapitel~\ref{chap:maf-realisierung}/\ref{chap:evaluation} \\
\bottomrule
\end{tabularx}
\end{table}


Über die Zyklen hinweg wurde vor allem die Aufgabenverteilung
ausdifferenziert. Modelle übernahmen semantisch offene Aufgaben.
Formal beschreibbare Prüfungen wurden deterministisch ausgeführt,
Verarbeitungsschritte durch explizite Workflows verbunden und
die fachliche Fortschreibung an menschliche Entscheidungspunkte gebunden.
Die zentrale Entwurfsfrage war damit, welche Verantwortung ein Agent
übernehmen sollte und wie sein Beitrag in einen kontrollierten
Gesamtablauf eingebettet werden konnte.

Die frühen Beobachtungen lassen sich nur unter Berücksichtigung der
jeweiligen Modellwahl einordnen. Nach wiederholt ausbleibenden
Werkzeugaufrufen mit lokal betriebenen Modellen kamen gehostete
Modelle wie GPT-4.1-mini zum Einsatz. Für qualitätsbezogene
Ledger-Versuche wurde Anfang Juli GPT-5.4 bevorzugt. Die
Entwicklungsnotizen begründen diese Wahl mit beobachteten Auslassungen
von Quellinhalten und zu nachsichtigen inhaltlichen Prüfurteilen
kleinerer Modelle. Kleinere Modelle blieben für kostengünstige
Funktionstests und gezielte Vergleiche im Einsatz; es gab keinen
vollständigen Wechsel zu nur einem Modell.

Die Modellwahl war damit Teil der explorativen Weiterentwicklung.
Da sich sowohl Modelle als auch Systemaufbau änderten, waren
beobachtete Unterschiede zunächst nicht eindeutig dem Design
zuzuordnen.

In ausgewählten explorativen Vergleichen blieb das Erzeugermodell
innerhalb der jeweiligen Vergleichsanordnung gleich. In einem dieser Vergleiche
wurde zusätzlich die direkte Erzeugung aus dem Transkript unter
denselben Modellbedingungen erneut ausgeführt. Damit entfiel ein
Modellwechsel als unmittelbare Erklärung für Unterschiede zwischen
den Varianten. Die begrenzte Vergleichsbasis erlaubte dennoch keine
kausale Zuordnung eines Effekts zum Design. Die Zyklen dokumentieren
die Entwicklungsrichtung; die Evaluation ausgewählter Mechanismen
folgt unter den in Kapitel~\ref{chap:evaluation} ausgewiesenen
Bedingungen.

% Historische Inline-Herkunftsnotizen vor M16; aktueller Befund im M16-Bericht.
% [Option A′ 01.09.: Mapping erweitert (Z7 → 4.6+4.7, Z8 → 4.8+4.9);
%  Z-Nummern bewusst UNVERÄNDERT (Kollegen-Entscheid).]
% [Einordnung: Kollegen-Vorschlag 01.09. („dosierte Agentik"-Runde);
%  deckt sich mit dem Projekt-Kernprinzip (Agent-Knoten in det.
%  Workflows mit Gates/Repair/Human-Autorisierung) und der 2.2/2.4.1-
%  Begriffsbasis (Agent-in-Workflow bleibt Agent).]
% [Kausalitäts-Check 01.09. (Kollegen-Nachfrage): die Verbindung IST
%  notes-belegt — phase01-notes Iteration 8 („Modell-/Provider-Wechsel:
%  OpenRouter als Alternative") nennt als Ausgangslage wörtlich die
%  Fehlerbilder der lokalen Modelle inkl. „erzeugten keine echten
%  fs_write-Toolcalls" → Nebeneinanderstellung bleibt, keine
%  Abschwächung nötig.]
% [M1 verifiziert 01.09. (E4-Entscheid: defensiv, keine Modellwertung):
%  phase01-notes Z. 432–438 (lokale qwen2.5-Modelle „instabil …
%  erzeugten keine echten fs_write-Toolcalls … Pseudo-Toolcalls"),
%  Z. 429 ff. (Iteration 8: OpenRouter-Wechsel, openai/gpt-4.1-mini);
%  „zuletzt gpt-5.4" (Autor-X) bewusst OHNE Modellnamen — Kap. 7 nennt
%  die Evaluations-Konfiguration.]
% [Vergleichsarm konkretisiert 01.09. (Kollegen-Punkt 5, verifiziert):
%  A'-Arm des evidence-first-Spikes = frisch erzeugtes freies Artefakt
%  unter gpt-5.4 (Run 20260630_145558_68149f, Commit c34b9a0), damit
%  B-vs-A nicht am alten Artefakt-/Modellstand hängt.]
% P10-6: Zyklenübersicht vor Beginn der Vertiefungen abschließen.
\clearpage
